% Fuentes: eq:Eab-generalizado, eq:Ephi-generalizado, eq:theta-generalizado
% y Ecuación métrica como extensión de L(g,R).
\begin{frame}[label=nav-frame-079]{Variación: caso geométrico y extensión escalar}
\TablaSetup
\footnotesize
\begin{tabular}{@{}p{.17\textwidth}p{.37\textwidth}p{.39\textwidth}@{}}
\TablaCabecera{Resultado} & \TablaCabecera{$L(g,R)$} & \TablaCabecera{$L(g,R,\phi,\nabla\phi)$} \\
\TablaLinea
$E_{ab}$ & $\begin{aligned}[t]&\mathcal R_{ab}-\tfrac12g_{ab}L\\&-2\nabla^m\nabla^nP_{amnb}\end{aligned}$ & $\begin{aligned}[t]&\mathcal R_{ab}-\tfrac12g_{ab}L\\&-2\nabla^m\nabla^nP_{amnb}\\&+\tfrac12\Pi_a u_b\end{aligned}$ \\[3mm]
$E_\phi$ & Sin ecuación escalar. & $F_\phi-\nabla_a\Pi^a$ \\[3mm]
Frontera & $\delta v^a$ & $\Theta^a=\delta v^a+\Pi^a\delta\phi$ \\[3mm]
Variaciones\newline de volumen & $E_{ab}\delta g^{ab}$ & $E_{ab}\delta g^{ab}+E_\phi\delta\phi$ \\
\TablaLinea
\end{tabular}\par
\vspace{4mm}
\small
$u_a=\nabla_a\phi$. En la columna escalar, $L$, $P^{abcd}$ y
$\mathcal R_{ab}$ corresponden a la teoría completa.
\end{frame}
